\chapter{SUSUNAN ORGANISASI TIM PENELITI}

{\scriptsize
\begin{longtable}{|c|p{2.8cm}|p{1.8cm}|p{1.9cm}|p{1.6cm}|p{3.5cm}|p{1.2cm}|}
\hline
\textbf{No} & \textbf{Nama Lengkap} & \textbf{NIP/NIDN/NIM} & \textbf{Kedudukan/ Jabatan} & \textbf{Rumpun Ilmu} & \textbf{Tanggung Jawab Kerja} & \textbf{Alokasi Waktu (jam/ minggu)} \\
\hline
\endfirsthead

\multicolumn{7}{c}{\textit{Lanjutan tabel}} \\
\hline
\textbf{No} & \textbf{Nama Lengkap} & \textbf{NIP/NIDN/NIM} & \textbf{Kedudukan/ Jabatan} & \textbf{Rumpun Ilmu} & \textbf{Tanggung Jawab Kerja} & \textbf{Alokasi Waktu (jam/ minggu)} \\
\hline
\endhead

\hline
\multicolumn{7}{r}{\textit{Bersambung...}} \\
\endfoot

\hline
\endlastfoot

1 & \ketuaNama & NIDN: \ketuaNIDN & Ketua/ Asisten Ahli & Rekayasa Perangkat Lunak &
\begin{itemize}[leftmargin=*, noitemsep, topsep=0pt]
\item Koordinasi tim penelitian
\item Desain arsitektur sistem
\item Pengembangan algoritma test case generation
\item Penyusunan laporan penelitian
\end{itemize} & 12 \\
\hline

2 & \anggotaSatuNama & NIDN: \anggotaSatuNIDN & Anggota/ Asisten Ahli & Basis Data dan Sistem Informasi &
\begin{itemize}[leftmargin=*, noitemsep, topsep=0pt]
\item Desain dan implementasi database
\item Integrasi dataset internal
\item Pengujian dengan aplikasi studi kasus (Laravel Blade, Inertia-React)
\item Analisis hasil pengujian
\end{itemize} & 10 \\
\hline

3 & \anggotaDuaNama & NIDN: \anggotaDuaNIDN & Anggota/ Lektor & Kecerdasan Buatan dan Data Science &
\begin{itemize}[leftmargin=*, noitemsep, topsep=0pt]
\item Validasi metode pengujian
\item Analisis metrik kinerja
\item Evaluasi cakupan pengujian
\item Penulisan artikel publikasi
\end{itemize} & 10 \\
\hline

4 & \anggotaTigaNama & NIDN: \anggotaTigaNIDN & Anggota/ Lektor & Jaringan Komputer dan Sistem Terdistribusi &
\begin{itemize}[leftmargin=*, noitemsep, topsep=0pt]
\item Implementasi web crawling
\item Integrasi Playwright framework
\item Pengujian cross-browser
\item Dokumentasi teknis
\end{itemize} & 10 \\
\hline

5 & \mahasiswaSatuNama & NIM: \mahasiswaSatuNIM & Mahasiswa & Teknik Informatika &
\begin{itemize}[leftmargin=*, noitemsep, topsep=0pt]
\item Pengumpulan dataset aplikasi
\item Implementasi fitur crawling
\item Testing prototipe awal
\item Dokumentasi kode program
\end{itemize} & 8 \\
\hline

6 & \mahasiswaDuaNama & NIM: \mahasiswaDuaNIM & Mahasiswa & Teknik Informatika &
\begin{itemize}[leftmargin=*, noitemsep, topsep=0pt]
\item Ekstraksi elemen HTML
\item Pengembangan modul parsing
\item Testing fitur ekstraksi
\item Pembuatan user manual
\end{itemize} & 8 \\
\hline

7 & \mahasiswaTigaNama & NIM: \mahasiswaTigaNIM & Mahasiswa & Teknik Informatika &
\begin{itemize}[leftmargin=*, noitemsep, topsep=0pt]
\item Implementasi test case generator
\item Pengujian dengan dataset
\item Perbaikan bug dan error
\item Pengumpulan data hasil testing
\end{itemize} & 8 \\
\hline

8 & \mahasiswaEmpatNama & NIM: \mahasiswaEmpatNIM & Mahasiswa & Teknik Informatika &
\begin{itemize}[leftmargin=*, noitemsep, topsep=0pt]
\item Integrasi dengan Playwright
\item Pengujian lintas browser
\item Analisis laporan hasil testing
\item Persiapan materi publikasi
\end{itemize} & 8 \\
\hline

\end{longtable}
}
